\documentclass[11pt,a4paper]{article}

\usepackage[margin=2.5cm]{geometry}
\usepackage[T1]{fontenc}
\usepackage{graphicx}
\usepackage{booktabs}
\usepackage{array}
\usepackage{amsmath}
\usepackage{xcolor}
\usepackage{caption}
\usepackage{float}
\usepackage{enumitem}
\usepackage[hidelinks]{hyperref}

\captionsetup{font=small,labelfont=bf}
\setlength{\parskip}{0.5em}
\setlength{\parindent}{0pt}

% ---------------------------------------------------------------------------
% Numbers are generated by data_scripts/*.py into stats.tex.
% If the scripts have not been run yet, fall back to a visible placeholder.
% ---------------------------------------------------------------------------
\InputIfFileExists{stats.tex}{}{}
\newcommand{\tbd}{\textbf{\textcolor{red}{[run scripts]}}}
\providecommand{\statTotalRaw}{\tbd}
\providecommand{\statCorrupt}{\tbd}
\providecommand{\statBadFormat}{\tbd}
\providecommand{\statConflictLabel}{\tbd}
\providecommand{\statExactDup}{\tbd}
\providecommand{\statDegenerate}{\tbd}
\providecommand{\statNearDupFlagged}{\tbd}
\providecommand{\statOutlierFlagged}{\tbd}
\providecommand{\statTotalClean}{\tbd}
\providecommand{\statTrain}{\tbd}
\providecommand{\statVal}{\tbd}
\providecommand{\statTest}{\tbd}
\providecommand{\statDatasetMean}{\tbd}
\providecommand{\statDatasetStd}{\tbd}

% Statements about the data that must be confirmed against the real figures
% are shown in red. Delete the \verify{...} wrapper once confirmed.
\newcommand{\verify}[1]{\textcolor{red}{#1}}

% Include a figure if it exists, otherwise show a placeholder box.
\newcommand{\safefig}[2][width=\linewidth]{%
  \IfFileExists{../figures/#2}%
    {\includegraphics[#1]{../figures/#2}}%
    {\fbox{\parbox{0.9\linewidth}{\centering Missing figure: \detokenize{#2}\\ run \texttt{data\_scripts/eda.py}}}}}

\title{\textbf{GeoScope: Explainable Land-Use Classification of Satellite Imagery}\\[0.3em]
\large Deliverable 1 --- Problem Definition and Data Acquisition}
\author{Ibrahim \textcolor{red}{[Surname, ID]} \and \textcolor{red}{[Teammate name, ID]}\\[0.5em]
Track A: Application Development\\
\small Repository: \url{https://github.com/Ibrahim392003/GeoScope}}
\date{October 2026}

\begin{document}
\maketitle

\begin{abstract}
GeoScope is a web application that classifies $64\times64$ Sentinel-2 satellite patches into ten
land-use classes and explains every prediction with a Grad-CAM heatmap. This report covers the first
two project phases. Phase~1 defines the problem, the target audience and the case for deep learning
over rule-based programming. Phase~2 documents the acquisition of the EuroSAT dataset, an
image-specific cleaning pipeline (corrupt files, duplicates, outliers), the feature strategy
(scaling, normalization, encoding, stratified splitting) and an initial exploratory data analysis.
All numbers in this report are generated directly by the scripts in \texttt{data\_scripts/}.
\end{abstract}

% =============================================================================
\section{Phase 1: Problem Definition}
% =============================================================================

\subsection{Problem Statement}
Knowing how land is used --- farmland, forest, housing, industry, water --- underpins urban
planning, agricultural monitoring, flood and disaster response, and environmental policy. Earth
observation satellites such as Copernicus Sentinel-2 now capture the whole of Europe every few days,
but turning that stream of pixels into reliable land-use maps remains a bottleneck. Manual
photo-interpretation needs trained experts and is slow and costly; as a result, official products such
as the CORINE Land Cover inventory are refreshed only about every six years.

Automatic classifiers based on deep learning can close this gap, but they bring a second problem:
they are \emph{black boxes}. A planner or analyst who receives a label such as ``Industrial'' with no
justification cannot tell whether the model looked at factory roofs or at an irrelevant artefact,
and therefore cannot decide how far to trust it. In addition, practitioners lack a simple way to
compare architectures (a small CNN, a modern efficient CNN and a Vision Transformer) on the same
data under the same conditions.

GeoScope addresses both pain points. It (i) classifies an uploaded satellite patch into one of ten
land-use classes in real time, (ii) overlays a Grad-CAM heatmap that highlights the image regions
that drove the decision, and (iii) benchmarks ResNet-18 \cite{he2016}, EfficientNet \cite{tan2019} and ViT \cite{dosovitskiy2021} so users can choose
the accuracy/speed trade-off that suits them.

\subsection{Target Audience}
\begin{itemize}[nosep]
  \item \textbf{Urban and regional planners} who need quick land-use checks and want to see
        \emph{why} a patch was labelled a certain way before acting on it.
  \item \textbf{Agricultural and environmental analysts} monitoring cropland, pasture and forest
        change who cannot afford manual interpretation of large areas.
  \item \textbf{GIS students and remote-sensing researchers} who want a transparent, reproducible
        benchmark of CNN and transformer models on a standard dataset.
\end{itemize}
The research gap addressed is the lack of an accessible tool that combines multi-architecture
benchmarking with built-in visual explanations for land-use classification. The scope of this
project is patch-level classification of RGB imagery; full-scene segmentation is out of scope.

\subsection{Why Machine Learning Instead of Rule-Based Programming?}
A rule-based system would encode expert knowledge as explicit conditions, for example
``if the patch is mostly green and smooth, then Forest''. This breaks down for land-use
classification for the following technical reasons.

\begin{enumerate}[nosep]
  \item \textbf{High intra-class variability.} The same class looks very different across seasons,
        crop stages, lighting conditions and countries. EuroSAT covers 34 European countries, so
        a threshold tuned in one region does not transfer to another.
  \item \textbf{Low inter-class separability in colour.} Several class pairs are visually close:
        \textit{Pasture} vs.\ \textit{HerbaceousVegetation}, \textit{AnnualCrop} vs.\
        \textit{PermanentCrop}, and \textit{Highway} vs.\ \textit{River} (long thin linear
        structures). Separating them depends on texture, geometry and spatial context, not on
        mean colour.
  \item \textbf{Spectral shortcuts are unavailable.} Classical remote-sensing rules rely on
        indices such as NDVI, which need the near-infrared band. GeoScope works on the RGB version
        of the data, where such indices cannot be computed.
  \item \textbf{Features cannot be enumerated by hand.} Deep networks learn a hierarchy of features
        directly from data: edges and colour blobs in early layers, textures such as field stripes
        and roof patterns in the middle layers, and object layouts such as road networks or
        building grids in deeper layers. Writing such detectors manually for ten classes is
        impractical.
  \item \textbf{Transfer learning reduces data needs.} Backbones pretrained on ImageNet can be
        fine-tuned on a moderately sized labelled set and reach high accuracy; the original EuroSAT
        benchmark reports about 98.6\% overall accuracy with a fine-tuned ResNet-50 \cite{helber2019}.
        No hand-crafted rule set has shown comparable performance on this task.
\end{enumerate}

Table~\ref{tab:rule-vs-ml} summarizes the comparison. A simple colour-histogram rule baseline
will still be implemented in Phase~3 as a lower bound so the benefit of deep learning is measured
rather than assumed.

\begin{table}[H]
\centering
\small
\begin{tabular}{@{}p{3.2cm}p{5.6cm}p{5.6cm}@{}}
\toprule
\textbf{Aspect} & \textbf{Rule-based} & \textbf{ML / DL (GeoScope)}\\
\midrule
Feature design & Manual thresholds on colour/texture & Learned hierarchically from pixels\\
Generalization & Fragile across season and region & Improves with diverse training data\\
Similar classes & Cannot separate on colour alone & Uses texture, shape and context\\
Maintenance & Rules rewritten for new regions & Fine-tune on new labelled data\\
Transparency & Inherently readable & Needs explanation (Grad-CAM \cite{selvaraju2017})\\
\bottomrule
\end{tabular}
\caption{Rule-based programming versus deep learning for land-use classification.}
\label{tab:rule-vs-ml}
\end{table}

The one advantage of rules --- transparency --- is exactly what GeoScope restores with Grad-CAM:
gradients flowing into a late convolutional layer are used to produce a coarse map of the regions
most responsible for the predicted class, letting a user check that the model is attending to
plausible evidence.

% =============================================================================
\section{Phase 2: Data Acquisition}
% =============================================================================

\subsection{Dataset}
GeoScope uses \textbf{EuroSAT} \cite{helber2019}: 27{,}000 labelled and geo-referenced Sentinel-2
patches of $64\times64$ pixels (10\,m ground resolution) in ten classes with 2{,}000--3{,}000 images
each (Table~\ref{tab:classes}). We use the RGB version distributed through \texttt{torchvision}.

\begin{table}[H]
\centering
\small
\begin{tabular}{@{}lrlr@{}}
\toprule
\textbf{Class} & \textbf{Images} & \textbf{Class} & \textbf{Images}\\
\midrule
AnnualCrop & 3{,}000 & PermanentCrop & 2{,}500\\
Forest & 3{,}000 & Residential & 3{,}000\\
HerbaceousVegetation & 3{,}000 & River & 2{,}500\\
Highway & 2{,}500 & SeaLake & 3{,}000\\
Industrial & 2{,}500 & Pasture & 2{,}000\\
\bottomrule
\end{tabular}
\caption{Official EuroSAT class sizes (total 27{,}000).}
\label{tab:classes}
\end{table}

\subsection{Collection}
The dataset is downloaded with \texttt{data\_scripts/download\_data.py}, which calls
\texttt{torchvision.datasets.EuroSAT(download=True)} and writes a log
(\texttt{report/download\_log.json}) with the timestamp, library version, image count and class list as
evidence of acquisition. After download, \statTotalRaw{} image files were found on disk. The
raw data is excluded from version control through \texttt{.gitignore} (\texttt{data/}), so anyone can
reproduce the collection with a single command.

\subsection{Cleaning Implementation}
Because the data are images rather than tables, the usual checks for missing values, duplicates and
outliers are translated into image-appropriate equivalents (Table~\ref{tab:cleaning}). The
implementation in \texttt{data\_scripts/clean\_data.py} never alters raw files; it writes a manifest
recording a status and a keep/discard decision per image, so every decision is auditable and
reproducible.

\begin{table}[H]
\centering
\small
\begin{tabular}{@{}p{3.3cm}p{4.6cm}p{3.4cm}r@{}}
\toprule
\textbf{Issue} & \textbf{Detection} & \textbf{Action} & \textbf{Count}\\
\midrule
Corrupt / unreadable & PIL \texttt{verify()} and decode & Removed & \statCorrupt\\
Wrong mode or size & Not $64\times64$ RGB & Removed & \statBadFormat\\
Label conflict & Identical pixels under two labels & All copies removed & \statConflictLabel\\
Exact duplicates & MD5 of decoded pixels & Keep first, drop rest & \statExactDup\\
Blank no-data patch & Contrast $<1$ and pure black/white & Removed & \statDegenerate\\
Near-duplicates & 64-bit dHash, Hamming $\le 3$ & Flagged only & \statNearDupFlagged\\
Statistical outliers & $3\times$IQR per class & Flagged only & \statOutlierFlagged\\
\midrule
\multicolumn{3}{@{}l}{\textbf{Images kept after cleaning}} & \textbf{\statTotalClean}\\
\bottomrule
\end{tabular}
\caption{Cleaning checks and their results on EuroSAT.}
\label{tab:cleaning}
\end{table}

\paragraph{Results.} Running the pipeline on all \statTotalRaw{} images found \statCorrupt{} corrupt files,
\statBadFormat{} images of the wrong size or mode, \statConflictLabel{} label conflicts, \statExactDup{} exact duplicates
and \statDegenerate{} blank patches, so no image had to be removed: EuroSAT is a carefully curated
dataset and the cleaning stage serves mainly as verification. The only items flagged were
\statNearDupFlagged{} near-duplicate images and \statOutlierFlagged{} statistical outliers, both of which are
kept for the reasons explained below. All \statTotalClean{} images proceed to modelling.

\paragraph{Missing values.} Image data has no empty cells; missing information appears as files that
cannot be decoded, images with the wrong size or colour mode, or class folders without images
(missing labels). The script detects all three, as well as any NaN in the extracted image features.

\paragraph{Duplicates.} Exact duplicates are found by hashing decoded pixel arrays (so differences in
file name or compression do not matter). If the same pixels appear under two different labels, the
label is unreliable and all copies are dropped. Near-duplicates are detected with a perceptual
difference hash; since satellite patches of uniform terrain (e.g.\ open water) can look alike
without being copies, these are \emph{flagged but kept}, and their share is reported rather than
silently removed. Duplicates that cross the train/test boundary would inflate test accuracy, which
is why exact duplicates are removed \emph{before} splitting.

\paragraph{Outliers.} Blank no-data patches (flat and pure black or white) carry no information and are removed. Uniform patches of normal brightness, such as calm water or dense forest, are \emph{not} removed: they are legitimate samples that the deployed app will also encounter, and deleting them would bias the dataset against those classes.
For the remaining images we compute brightness, contrast and the variance of the Laplacian
(sharpness) and flag values outside $3\times$IQR of the image's own class
(Figure~\ref{fig:outliers}). Statistical outliers in satellite imagery are usually legitimate scene
variation (shadows, snow, cloud edges) and teach the model robustness, so they are flagged rather than deleted;
the script offers \texttt{--drop-outliers} if later experiments show they hurt accuracy.

\begin{figure}[H]
\centering
\begin{minipage}{0.49\linewidth}\safefig{cleaning_summary.png}\end{minipage}
\hfill
\begin{minipage}{0.49\linewidth}\safefig{outlier_detection.png}\end{minipage}
\caption{Left: images removed by the cleaning pipeline and why. Right: brightness vs.\ contrast
with per-class outliers highlighted in red.}
\label{fig:outliers}
\end{figure}

\subsection{Feature Strategy}
Implemented in \texttt{data\_scripts/preprocess.py}.

\begin{description}[style=nextline,leftmargin=1.2em]
  \item[Resizing.] Patches are upscaled from $64\times64$ to $224\times224$, the input size on
        which ResNet-18, EfficientNet and ViT-B/16 were pretrained.
  \item[Scaling.] \texttt{ToTensor} maps 8-bit pixel values from $[0,255]$ to $[0,1]$.
  \item[Normalization.] Each channel is standardized, $x' = (x-\mu)/\sigma$. For pretrained backbones
        we use the ImageNet statistics $\mu=(0.485,0.456,0.406)$, $\sigma=(0.229,0.224,0.225)$, because
        the pretrained weights expect inputs distributed this way. We also computed the EuroSAT
        training-set statistics ($\mu=$ \statDatasetMean, $\sigma=$ \statDatasetStd), which will be
        used if a model is trained from scratch. Figure~\ref{fig:norm} shows the effect.
  \item[Encoding.] The ten class names are integer-encoded in alphabetical order
        (\texttt{AnnualCrop}=0 \ldots \texttt{SeaLake}=9), saved in \texttt{label\_map.json}. Cross-entropy
        loss takes integer targets directly, so one-hot encoding is not needed.
  \item[Splitting.] A stratified $70/15/15$ train/validation/test split with fixed seed 42 gives
        \statTrain{} / \statVal{} / \statTest{} images and keeps class proportions identical in
        every split.
  \item[Augmentation (train only).] Random horizontal and vertical flips; land-use patches have no
        canonical orientation, so flips create realistic new views.
  \item[Class imbalance.] The largest class is 1.5$\times$ the size of the smallest (3{,}000 vs.\ 2{,}000
        before cleaning). This mild imbalance does not justify oversampling or undersampling; we will
        instead report macro-averaged F1 alongside accuracy and keep class-weighted loss as an option.
\end{description}

\begin{figure}[H]
\centering
\safefig[width=0.9\linewidth]{normalization_effect.png}
\caption{Per-channel pixel distributions after scaling to $[0,1]$ (left) and after ImageNet
normalization (right).}
\label{fig:norm}
\end{figure}

\subsection{Exploratory Data Analysis}
All charts are produced by \texttt{data\_scripts/eda.py}.

\paragraph{Class distribution.} Figure~\ref{fig:dist} shows the raw and cleaned class counts. The dataset
is nearly balanced (largest class 1.5$\times$ the smallest). Because the cleaning stage removed no images
(Table~\ref{tab:cleaning}), the raw and cleaned bars are identical.

\begin{figure}[H]
\centering
\safefig{class_distribution.png}
\caption{Number of images per class before and after cleaning.}
\label{fig:dist}
\end{figure}

\paragraph{Visual inspection.} Figure~\ref{fig:grid} shows random samples from each class. Classes are
defined by spatial structure rather than colour: Industrial shows bright rooftops, Residential a fine grid
of houses, Highway and River long linear features, while Forest and SeaLake are nearly featureless.
Appearance also varies within a class (SeaLake ranges from teal to navy), and similar pairs such as
AnnualCrop and PermanentCrop are easy to confuse. This supports the argument that fixed colour rules are
insufficient.

\begin{figure}[H]
\centering
\safefig[height=0.62\textheight]{sample_grid.png}
\caption{Random samples from each of the ten classes.}
\label{fig:grid}
\end{figure}

\paragraph{Pixel intensities and class statistics.} Figure~\ref{fig:hist} shows the global channel
histograms and Figure~\ref{fig:bright} the per-class brightness and contrast.
SeaLake and Forest are the darkest classes (mean brightness 67 and 58 on a 0--255 scale) and have
very low contrast (below 4), reflecting their uniform appearance. Industrial patches have the highest
contrast (41.7) and sharpness, followed by Residential, which fits the dense mix of roofs, roads and shadows in built-up areas.
Brightness and contrast therefore separate some classes (e.g.\ SeaLake vs.\ Industrial) but overlap
strongly for others (e.g.\ AnnualCrop, PermanentCrop and Industrial have similar mean brightness of
118--126), so simple intensity cues carry some signal but are not sufficient.

\begin{figure}[H]
\centering
\safefig[width=0.7\linewidth]{channel_histograms.png}
\caption{Pixel intensity distribution per RGB channel (100 sampled images per class).}
\label{fig:hist}
\end{figure}

\begin{figure}[H]
\centering
\safefig{brightness_by_class.png}
\caption{Brightness and contrast by class.}
\label{fig:bright}
\end{figure}

\paragraph{Feature correlation.} Figure~\ref{fig:corr} shows the correlation between image-level
statistics. The mean R, G and B values are strongly correlated with each other and with
brightness, meaning colour alone is largely redundant information, which further motivates models
that learn spatial structure instead.

\begin{figure}[H]
\centering
\safefig[width=0.55\linewidth]{feature_correlation.png}
\caption{Correlation matrix of image-level features.}
\label{fig:corr}
\end{figure}

% =============================================================================
\section{Summary and Next Steps}
% =============================================================================
Phase~1 established that automatic, \emph{explainable} land-use classification is a real need and that
deep learning is justified over rule-based programming because of the high variability and low
colour separability of the classes. Phase~2 delivered a reproducible pipeline: scripted collection,
an auditable cleaning stage, a stratified split, a documented normalization strategy and an initial
EDA, leaving \statTotalClean{} clean images for modelling.

Planned work for the next phases:
\begin{itemize}[nosep]
  \item Train and benchmark ResNet-18, EfficientNet and ViT-B/16 with transfer learning, plus a
        colour-histogram rule-based baseline; report accuracy, macro-F1, confusion matrices and
        inference speed.
  \item Implement Grad-CAM for the CNNs. \emph{Risk:} ViT has no convolutional feature maps, so
        explanations will need an attention-based variant (e.g.\ attention rollout or Grad-CAM on
        token embeddings); this will be evaluated early.
  \item Check for near-duplicate leakage across splits and tune the outlier policy using validation
        results.
  \item Build and deploy the web interface.
\end{itemize}

\begin{thebibliography}{9}
\bibitem{helber2019} P.~Helber, B.~Bischke, A.~Dengel, D.~Borth.
  EuroSAT: A Novel Dataset and Deep Learning Benchmark for Land Use and Land Cover Classification.
  \emph{IEEE Journal of Selected Topics in Applied Earth Observations and Remote Sensing}, 2019.
\bibitem{selvaraju2017} R.~R.~Selvaraju et al.
  Grad-CAM: Visual Explanations from Deep Networks via Gradient-based Localization.
  \emph{ICCV}, 2017.
\bibitem{he2016} K.~He, X.~Zhang, S.~Ren, J.~Sun. Deep Residual Learning for Image Recognition. \emph{CVPR}, 2016.
\bibitem{tan2019} M.~Tan, Q.~V.~Le. EfficientNet: Rethinking Model Scaling for Convolutional Neural Networks. \emph{ICML}, 2019.
\bibitem{dosovitskiy2021} A.~Dosovitskiy et al. An Image is Worth 16x16 Words: Transformers for Image Recognition at Scale. \emph{ICLR}, 2021.
\end{thebibliography}

\end{document}